\section{本讲主线：模型训练结束，系统问题才刚开始}

CS336 的大部分课程从数据、优化、并行训练和 post-training 解释“模型是怎样得到的”。Dan Fu 的嘉宾课把视角翻到另一面：当权重文件已经存在，怎样把它变成一个能承受真实流量、满足延迟目标、持续输出 token 的服务？更重要的是，推理系统暴露出的瓶颈会不会反过来改变 kernel、并行策略，甚至模型架构本身？

这堂课可以压缩成一条因果链：\term{workload shape} 决定 SLA，SLA 决定 scheduler 和 cache 策略，scheduler 决定 GPU 上出现怎样的空洞，kernel 级空洞又决定哪些架构值得设计。优化对象不再是一个孤立的 forward pass，而是端到端的“请求到 token”系统。

\begin{importantbox}{全栈推理的统一目标}
不要单独最大化某一项 FLOPs、带宽利用率或 batch size。真正的目标是在给定 workload 分布和 tail-latency 约束下，最大化可持续吞吐，同时保留模型质量、故障恢复能力和工程可维护性。
\end{importantbox}

\begin{figure}[H]
\centering
\includegraphics[width=0.92\textwidth]{00-08-40-inference-engine.jpg}
\caption{推理引擎的高层地图：请求进入后依次经过调度、缓存、执行、并行与输出。}
\figsource{课堂视频 00:08:10--00:09:38；该页由讲者用生成式工具制作，只取高层结构，不采信细小文字。}
\end{figure}

\subsection{怎样阅读本讲的图}

讲者在 00:07:57 明确说明，前半部分若干幻灯片由 Nano Banana Pro 生成，“高层正确，但细看文字可能错误”。因此，本讲义只把这些图当作架构示意：组件名称、指标定义、公式和数字均以课堂口述、Parcae 论文、Together AI 的 CPD 说明和 HazyResearch 的 megakernel 资料交叉核验。

\begin{warningbox}{图像不是事实来源}
生成式幻灯片尤其容易在坐标轴、单位、缩写和小标签上产生幻觉。读图时应先从讲者口述建立机制，再用图确认组件关系；不能从模糊小字反向推导结论。
\end{warningbox}

\subsection{术语表：本讲反复出现的对象}

为了避免系统术语在后文反复跳转，本节先把最容易混淆的对象放在同一张表里。读者应特别区分“延迟指标”和“执行阶段”，也要区分保存历史状态的 KV cache 与真正执行矩阵运算的 kernel；后文每次优化都可以回到这些对象判断它改变了什么。

\begin{table}[H]
\centering
\begin{tabular}{L{0.19\textwidth}L{0.29\textwidth}L{0.42\textwidth}}
\toprule
术语 & 简明定义 & 在本讲中的作用 \\
\midrule
Prefill & 并行处理输入上下文并建立 KV state & 主要影响 TTFT，长 prompt 成本高 \\
Decode & 自回归地产生后续 token & 主要影响 TBT，常受权重与 KV 搬运限制 \\
KV cache & 历史 token 的 attention key/value 状态 & 支撑 prefix reuse、多轮对话与长上下文 \\
Scheduler & 决定请求何时、在哪个 worker、与谁组成 batch & 平衡吞吐、显存、cache locality 与 tail SLA \\
Kernel & GPU 上执行某段算子的程序 & launch、同步、load 和 tail effect 决定底层效率 \\
Recurrence & 让 activation 重复经过同一参数 block & 增加计算而不同比例增加权重参数 \\
\bottomrule
\end{tabular}
\caption{本讲核心术语表。}
\end{table}

\subsection{本章小结}

本讲不是一份“推理技巧清单”，而是一个跨层设计案例：从流量到服务指标，从服务指标到缓存与调度，从 GPU 时间线到 megakernel，再从系统约束回到 Parcae 这类新架构。

